\documentclass[]{article}

%opening
\title{Computing Science PhD Research Proposal}
\author{Your Name}

\begin{document}

\maketitle

\section{Introduction}
\label{sec:intro}

Describe the area you want to work in (e.g.\ CHERI systems \cite{woodruff2014cheri,watson2020cheri})
and the precise topic you want to focus on (i.e.\ some kind of runtime
systems research).

Explain the structure of the rest of this proposal document.

\section{Hypothesis}
\label{sec:hyp}

State your initial research hypothesis here. This should be a falsifiable
assertion which you can investigate throughout your research project.

If required, you might break the investigation down into a number of
smaller \textit{research questions}.

\begin{itemize}
\item \textbf{RQ1}: first research question
\item \textbf{RQ2}: second research question
\item \textbf{RQ3}: third research question
\end{itemize}



\section{Background Survey}
\label{sec:bg}

You need to outline the motivating context, and cite relevant papers.
Use BibTeX to get the citations.

You should also mention interesting systems papers that link in
with your research hypothesis.

\section{Research Plan}
\label{sec:plan}

Here you describe in some detail what you would like to do in the next
year.



After you have settled on a research plan, you should plan out the work you
want to do into brief work packages and deliverables, then plot these on a
Gantt chart. There is a nice LaTeX package for drawing Gantt charts, called
pgfgantt; I highly recommend using it.


\bibliographystyle{abbrv}
\bibliography{proposal}

\end{document}
